\section{Introduction}
\label{sec:intro}

Most reinforcement learning for language models optimizes a single scalar
reward and returns a single policy. Many real objectives are not like that. In
attribute-controlled rewriting, a system must satisfy two objectives in tension
--- change the attribute, preserve the meaning --- and the useful artifact is
not one operating point but the whole tradeoff curve, so that a downstream user
can choose where on it to sit. Training one policy per point is wasteful and
gives no guarantee that the points are mutually consistent.

We study this in the prompt-optimization setting~\citep{deng2022rlprompt}, where the trainable object is
a short discrete prompt that steers a \emph{frozen} task language model. Our
policy is conditioned on a scalar $\lmb$ that indexes how tightly content
preservation is required, so a single trained policy traces the entire
content-versus-sentiment curve by varying its conditioning input at test time.
This makes exploration unusually load-bearing: a policy that collapses onto one
prompt is not merely suboptimal, it becomes $\lmb$-independent and collapses the
curve to a point.

That observation motivates the algorithm we study. GRPO~\citep{shao2024deepseekmath}, the default choice for
this kind of problem, estimates a group-relative advantage and ascends it inside
a clipped trust region. But in a pipeline that takes one gradient update per
batch of rollouts --- as prompt optimization naturally does, because rollouts
are expensive and the policy is small --- the importance ratio is identically
one in value and the clip is provably inert, leaving the KL penalty as the objective's
only regularizer, and leaving its unconstrained optimum a deterministic policy.
\rrebel{}, a robust extension of REBEL~\citep{gao2024rebel}, instead treats
improvement as regression: it fits the $\beta$-scaled
log-ratio difference of two sampled prompts to their observed reward difference,
under a \emph{robust} discrepancy, with rewards standardized within each group.
Why this works, and when it should, has so far been argued informally. We give
it a theory in the style of DPO~\citep{rafailov2023dpo}, which says what the
objective converges to and what each design choice buys. The theory also
predicts the regime where group standardization matters: one in which the
reward's scale varies enormously across the contexts a single policy must
serve. We measure that this multi-objective problem is squarely in that
regime.

\paragraph{Contributions.}
\begin{itemize}\itemsep2pt
  \item A $\lmb$-conditioned prompt policy for multi-objective controlled
        rewriting, trained against a constrained reward in which $\lmb$ sets a
        content floor (Section~\ref{sec:setup}).
  \item A theory of \rrebel{}-std (Section~\ref{sec:theory}). Its population
        minimizer is a KL-regularized optimum for any discrepancy and any
        full-support sampler, with group standardization turning one global
        temperature into a per-context one. Consequently every context gets the
        same trust-region radius, and a lagged reference performs mirror ascent
        with standardized steps. Pairwise differencing symmetrizes reward
        noise, so robust discrepancies do not bias the target. GRPO without a KL
        term is the special case of \rrebel{}-$\ell_2$-std with lag one. Every
        statement is proved and checked numerically against the implementation.
  \item A held-out test-split evaluation of the full tradeoff curve for two
        \rrebel{} and two GRPO variants (Section~\ref{sec:results}):
        $\vNsent$ sentences, $\vNlam$ content floors, repeated task-model
        seeds, and paired bootstrap intervals. \rrebel{} with a Huber
        discrepancy has the best curve (hypervolume $\vHuberHV$ against
        $\vGentHV$, reward $\vHuberRew$ against $\vGentRew$). We also report
        how well each policy honours its floor, what prompts it learned, how its
        curve evolved over training, and a direct measurement of the reward-scale
        variation that motivates standardization.
  \item Two negative results that bound the comparison. A GRPO baseline
        anchored to the pretrained backbone, with a genuine trust region, is the
        weakest arm. And the $\ell_1$ variant's curve regresses late in
        training, which development-set early stopping avoids.
\end{itemize}
